\section{谱聚类把连通性变成特征向量}
\label{sec:13-Machine-Learning/05-Clustering-and-Data-Geometry/02}

同一批用户数据，你把上一节的 K-means 跑完，画了张二维投影，结果看得人发愁：屏幕上明明是两条平行的弯带，算法却把它们从中间竖着劈开，每个簇里一半来自这条带一半来自那条带。你调 \(K\)、换标准化、多跑几次初始化，切法只是换个角度，从没沿着带子走。

问题不在优化，在提问方式。K-means 问的是``你离哪个中心近''，而一条弯带上首尾两端的点隔得很远，中间却密密麻麻连成一串。人眼判定它们同属一群，依据是``从这头能一路走到那头''，而不是``它们离某个点都近''。把``近''换成``连通''，需要的就不再是坐标和中心，而是一张图。

谱聚类的做法是：先把样本变成一张带权图，再找一种把图切开而尽量少剪断强边的分法。这个提问方式换掉之后，簇的定义也就跟着换了。

\subsection{相似图是你亲手造的第一件东西}

起点是一个对称非负的权重矩阵 \(W_{ij}=W_{ji}\ge 0\)，常见的造法是先建 \(k\) 近邻图，再在保留的边上放高斯权重

\[
W_{ij}=\exp\left(-\frac{\norm{x_i-x_j}^2}{2\sigma^2}\right).
\]

不在近邻表里的点对权重取零——没有边就没有直接影响。度矩阵汇总每个节点连出去的总权重，

\[
D=\diag(d_1,\ldots,d_n),\qquad d_i=\sum_j W_{ij},
\]

非归一化图 Laplacian 则定义为 \(L=D-W\)。

后面所有的谱分析都建在这个 \(W\) 上，所以先说清楚它有多脆。近邻数 \(k\) 取小了，本来连着的一簇会断成几截，图里出现数据中并不存在的裂缝；取大了，两个不同的簇之间会被架起捷径，图把不该连的连上了。\(\sigma\) 取小了只剩极局部的边，大量点变成孤岛；取大了所有权重趋于相等，图彻底失去分辨力。谱方法并没有取消尺度选择这件事，它把这个选择从``中心怎么算''搬到了``图怎么建''，而搬过去之后它反而更隐蔽，因为你不再直接看得到它。

\subsection{Laplacian 二次型给``被拆散''定价}

给每个节点赋一个实数，得到向量 \(f\in\R^n\)。展开可得

\[
f^{\trans}Lf=\frac12\sum_{i,j}W_{ij}(f_i-f_j)^2 .
\]

这个式子可以直接读：一条权重很大的边，若两端取值差很多，\(W_{ij}(f_i-f_j)^2\) 就很大，罚得重；强连接的节点取值相近则几乎不罚。取 \(f=\mathbf 1\)，所有差为零，于是 \(L\mathbf 1=0\)，常数向量总在零特征空间里。二次型非负说明 \(L\) 半正定，而 \(L\) 对称保证了特征向量可以取成一组正交基，这正是\hyperref[sec:03-Linear-Algebra/06-Symmetric-and-Positive-Definite-Matrices/01]{``对称矩阵与谱定理：什么时候特征方向一定正交''}给出的保证。

若图恰好有 \(c\) 个互不相连的连通分量，零特征值的重数就是 \(c\)：每个分量上取 1、其余取 0 的指示向量都让二次型为零，它们张成整个零空间。理想情形下连通分量直接就是簇，连特征分解都不必细看。

现实的图没这么干净，簇之间总有零星弱边把一切粘在一起，零特征值的重数掉回 1。但结构并没有消失，它挪到了紧挨着零的那几个小特征值上。对应的特征向量沿弱边可以变化，沿强边几乎不变——它们在描述图上``哪些地方允许被撕开''。于是聚类问题变成了：找几个最小的非零特征值方向。

\subsection{离散的割经过一次松弛变成特征问题}

先看两分的情形。用取值在 \(\pm1\) 的向量表示两侧归属，被切断的边权之和可以写成 \(f^{\trans}Lf\) 的常数倍。直接在 \(\set{\pm1}^n\) 上优化是组合问题，状态数随节点数指数增长。

松一步：允许 \(f\) 取实值，加上 \(f\perp\mathbf 1\) 去掉平凡的常数解，再加 \(\norm f=1\) 固定尺度，问题变成 Rayleigh 商的最小化

\[
\min_{f\perp\mathbf 1,\;\norm f=1} f^{\trans}Lf .
\]

按变分刻画，解是第二小特征值对应的特征向量，习惯上叫 Fiedler 向量。按它的符号切，或者沿某个阈值切，就得到一个近似的二分。多簇时取前 \(k\) 个最小特征向量，把节点 \(i\) 的新坐标定义为这 \(k\) 个向量的第 \(i\) 个分量拼成的行向量，再在这组谱坐标上跑一次 K-means。

\begin{center}
\begin{tikzpicture}[>=stealth]
% ---- 左：原空间，两个同心环 ----
\draw[thick] (0,0) rectangle (3.4,3.4);
\foreach \a in {0,15,...,345}
  \draw ({1.7+1.35*cos(\a)},{1.7+1.35*sin(\a)}) circle (2.0pt);
\foreach \a in {0,30,...,330}
  \fill ({1.7+0.5*cos(\a)},{1.7+0.5*sin(\a)}) circle (2.0pt);
\draw[very thick] (1.7,0.05) -- (1.7,3.35);
\node at (1.7,-0.5) {\footnotesize 原空间：K-means 竖着切};
% ---- 右：谱坐标上的一维展开 ----
\begin{scope}[shift={(4.8,0)}]
\draw[thick] (0,0) rectangle (3.4,3.4);
\draw[->] (0.2,1.7) -- (3.25,1.7);
\node[font=\footnotesize] at (1.55,1.98) {\(f_2\)};
\foreach \p in {(0.55,1.7),(0.62,1.92),(0.70,1.50),(0.78,1.82),(0.50,1.58),(0.85,1.70),(0.66,2.05),(0.74,1.38),(0.58,1.44),(0.80,2.00),(0.90,1.55),(0.46,1.84)}
  \fill \p circle (2.0pt);
\foreach \p in {(2.30,1.70),(2.42,1.95),(2.55,1.45),(2.68,1.85),(2.20,1.55),(2.80,1.72),(2.48,2.08),(2.60,1.35),(2.34,1.40),(2.72,2.02),(2.90,1.60),(2.16,1.88),(2.86,1.42),(2.26,2.05),(2.62,1.62),(2.38,1.62),(2.52,1.78),(2.76,1.50),(2.10,1.72),(2.94,1.86),(2.44,1.25),(2.66,2.15),(2.22,1.28),(2.84,2.12)}
  \draw \p circle (2.0pt);
\node at (1.7,-0.5) {\footnotesize 谱坐标：两环各自缩成一团};
\end{scope}
\end{tikzpicture}
\end{center}

\emph{同一批点，左边按坐标看无从下手，右边在 Fiedler 向量给的一维坐标上一刀就分开。}

图里左边是两个同心环。K-means 的两个中心必然落在环之间的空地上，边界是一条经过圆心的直线，把两个环同时劈开——它想找的团状结构根本不存在。右边是同一批点在 Fiedler 向量上的坐标：内环的点因为彼此强连通而被压到一个值附近，外环同理，两团之间空出一大段。最后那一步 K-means 并没有被抛弃，它只是被挪到了一个已经好切的空间里跑。谱聚类做的是换表示，分类的动作还是老一套。

松弛这一步值得留个心眼。连续解与最优的离散割之间没有一般性的等价，两者的差距可以被构造得很大，特征向量给出的只是一个有理论依据的候选，不是最优割的伪装。取多少个特征向量同样是个选择：常见的启发是看特征值序列 \(\lambda_1\le\lambda_2\le\cdots\) 中哪里出现明显跳跃，即所谓 eigengap，但和上一节的肘部法则一样，跳跃不明显时这个规则不给答案。

\subsection{归一化决定你在平衡什么}

只看被切断的边权总和会出问题。孤立一个度数很低的单点，切断的权重极小，代价函数很满意，聚类任务却一无所获。

Ratio cut 把代价除以每块的节点数，normalized cut 除以每块的总度数，两者都在要求切出来的块得足够大或足够密。反映到矩阵上，是两种归一化的 Laplacian：

\[
L_{\mathrm{sym}}=I-D^{-1/2}WD^{-1/2},
\qquad
L_{\mathrm{rw}}=I-D^{-1}W .
\]

后者有个直观读法：\(D^{-1}W\) 恰好是图上随机游走的转移矩阵，\(L_{\mathrm{rw}}\) 的小特征值对应游走中衰减最慢的模式，也就是概率质量长期滞留其中、很少跨出去的那些区域，这与\hyperref[sec:06-Random-Processes/03-Markov-Chains/01]{``Markov 性质、转移矩阵与多步演化''}里描述的多步演化是同一件事。簇因此获得了一个动力学定义：随机游走者短期内跑不出去的地方。

两种归一化背后是不同的平衡准则，在度数分布悬殊的图上可能给出不同分区，比如一个密集核心挂着一圈稀疏外围时。名字里都有``谱聚类''并不意味着公式可以互换，先确认你想平衡的是节点数还是总度数。

\subsection{谱嵌入好看是被构造出来的}

谱嵌入沿相似图的连通方向展开数据，所以它的二维图里出现干净的团块并不构成证据——这正是它被设计出来做的事。检验的办法是动参数：换一组 \(k\) 和 \(\sigma\) 重建图，若团块随之消散或重组，结论主要来自建图选择。样本密度不均匀时这一点尤其要紧，固定近邻数会在稠密区把一簇切碎，固定带宽会让稀疏区整片断开，而真实数据的密度很少是均匀的。

高维带来的麻烦和上一节一样，只是隐藏得更深。权重矩阵是从 \(\norm{x_i-x_j}\) 算出来的，距离一旦集中，所有 \(W_{ij}\) 就趋于相等，图退化成近似完全图，Laplacian 的小特征值挤成一堆，特征向量的方向由噪声决定。此时谱嵌入照样会画出漂亮的团块，因为它总能找到某个方向把点摊开。要让连通性假设站得住，数据得实际躺在某个低维流形上，邻域才既局部又连贯，这正是\hyperref[sec:13-Machine-Learning/04-Dimensionality-Reduction/04]{``流形假设与非线性降维''}讨论的前提。

计算上还有一层要单独交代。稠密的 \(W\) 要 \(O(n^2)\) 空间，完整特征分解更贵；稀疏近邻图配上 Lanczos 一类的迭代求解器可以把规模做上去，但近邻搜索的近似误差会直接写进图结构，随后被特征分解放大或掩盖。数值误差、建图误差和统计不确定性是三件事，一张好看的嵌入图无法同时回答它们。

回到本单元开头的处境。上一节的 K-means 假设簇围绕中心并且各向同性，本节的谱聚类假设簇内部由高权重路径连通、簇间连接稀薄。同一批点在两套假设下会得到不同的分组，而这两组结果都是各自准则下的正确解——没有一个隐藏的标准答案在背后等着被逼近。因此报告聚类结果时，``分成了几类''只是其中最不重要的一项；该一并交代的是你选了哪把尺子、哪张图、哪个 \(K\)，以及换一套之后还剩下多少结论站得住。
